\section{Failure Modes}
\label{sec:failure}

A score that is useful in deployment must be characterized where it errs as well
as where it succeeds. We examine both directions of error on the frontier cohort,
and because the venue is open we read the actual human reviews to explain each.
Table~\ref{tab:casestudy} (Appendix~\ref{app:cases}) places AIPR's stated weakness
beside a verbatim excerpt from the human reviews for a matched set of submissions
spanning the patterns below; we paraphrase here and quote in full there.

\subsection{What the score catches (and why)}
For rejected submissions the score rated low, the in-depth review's stated
weaknesses align closely with the reasons human reviewers gave---and at the level of
the specific defect, not merely the overall verdict. The recurring,
correctly-flagged patterns, each instantiated in Table~\ref{tab:casestudy}, are:
\begin{itemize}\itemsep1pt
  \item \textbf{Unquantified or synthetic data construction}: AIPR flags that the
  reported gains may hinge on how a synthetic benchmark was built, and the human
  reviewers independently interrogate the same construction.
  \item \textbf{Thin or missing baselines}: AIPR names the absent comparisons
  (e.g.\ ``only ERM and static BIRM''); reviewers name the same gap (``too few
  baselines'').
  \item \textbf{Overclaiming and poor situating}: AIPR flags a gap between the
  title/abstract claims and the evidence, and missing prior work; reviewers lead
  with the identical objection.
\end{itemize}
In each case the model and the reviewers point at the same table, claim, or
omission, which is what makes a low score a usable flag rather than a coincidence.

\subsection{Where the score is wrong}
We report both error classes explicitly, with examples in Table~\ref{tab:casestudy}.
\paragraph{False negatives: rejected work the score rated highly.}
The dangerous case for a triage tool is a weak paper scored high. In the observed
examples these concentrate where the rejection rests on a \emph{contribution-level
judgment} that a competent manuscript surface does not betray: reviewers call the
method ``an arguably trivial combination of existing'' techniques with ``little
conceptual novelty,'' while AIPR, scoring a well-written and well-executed paper,
places it in the upper band. A second, structural cause is a defect the system
cannot see at all---it reads the manuscript as written and cannot execute an
experiment, reproduce a result, or re-derive a proof---so where rejection hinges on
such a defect a high score is expected and is not evidence against the bounded claim
of \S\ref{sec:results}. Either way, the rate of these errors at the top of the score
range is quantified directly in the reliability curve (Fig.~\ref{fig:validation}b).

\paragraph{False positives: strong work the score underrated.}
The mirror error is accepted work scored low. In the observed examples AIPR's lower
score reflects a \emph{genuine} concern---an over-reaching theoretical claim, or a
theory whose assumptions do not match practice---that a human reviewer often raised
as well, but that the community forgave in light of the contribution: the papers
were accepted as orals on their merits. AIPR underweights a strong contribution when
it carries a real but pardonable gap. These cost a second human look, which is the
intended action anyway, and they bound any use of the score for positive selection.

\subsection{Reading of the failure modes}
The two error classes differ in their consequences for deployment. A false
negative at the \emph{top} of the score range (a weak paper scored high) is the
dangerous case for a triage tool, and we quantify its rate directly in the
reliability curve (Fig.~\ref{fig:validation}b). A false positive (a strong paper
scored low) costs a second human look, which is the intended action anyway. This
asymmetry is what makes the narrow claim, to flag weak work for human attention,
defensible as a human-reviewed triage signal under the studied conditions even
where the score ranks strong work poorly.
